% Superpapers paper skeleton — minimal LaTeX template for empirical research
% Compile with: ./skills/compile-latex/scripts/compile.sh paper/paper.tex
%
% Language: default English. For other languages, uncomment the babel line
% and adjust. All preamble comments remain in English.

\documentclass[12pt,a4paper]{article}

% --- Encoding and language ---
\usepackage[utf8]{inputenc}
\usepackage[T1]{fontenc}
% \usepackage[english]{babel}      % default
% \usepackage[brazilian]{babel}    % for pt-BR
% \usepackage[spanish]{babel}      % for es
% \usepackage[french]{babel}       % for fr

% --- Typography ---
\usepackage{lmodern}
\usepackage{microtype}
% \usepackage{fontspec}            % uncomment to use xelatex with custom fonts
                                   % (then remove inputenc, fontenc, and lmodern above)

% --- Math ---
\usepackage{amsmath}
\usepackage{amssymb}
\usepackage{amsthm}

% --- Tables and figures ---
\usepackage{booktabs}
\usepackage{threeparttable}
\usepackage{graphicx}
\usepackage{float}

% --- Bibliography (choose one) ---
% Author-year is the plugin default; numeric only if the target journal requires it.
% If switching to a journal class that loads natbib itself (e.g., elsarticle),
% set the mode via the class option instead:
%   \documentclass[preprint,12pt,authoryear]{elsarticle}
\usepackage[authoryear,round]{natbib}  % natbib style
% \usepackage[style=authoryear,backend=biber]{biblatex}  % biblatex alternative
% \addbibresource{references.bib}                        % for biblatex

% --- Links ---
\usepackage[colorlinks=true,allcolors=blue]{hyperref}

% --- Layout ---
\usepackage[margin=1in]{geometry}
\usepackage{setspace}
\onehalfspacing

% --- Metadata ---
\title{<Title of the Paper>}
\author{<Author Name>\thanks{<Affiliation and contact>}}
\date{\today}

\begin{document}

\maketitle

\begin{abstract}
\noindent <Abstract: 150-250 words. State the research question, data, method,
main finding, and contribution.>

\noindent\textbf{Keywords:} <keyword1, keyword2, keyword3>

\noindent\textbf{JEL codes:} <JEL1, JEL2>  % remove if not an economics paper
\end{abstract}

\section{Introduction}
\label{sec:intro}

<Motivation. Research question. Contribution. Preview of findings. Roadmap.>

\section{Conceptual Framework}
\label{sec:framework}

<Theoretical mechanisms and hypotheses: why should X affect Y? Candidate
channels the Results section will test. The literature review itself belongs
in the Introduction, not in a standalone section.>

\section{Data}
\label{sec:data}

<Sources, sample selection, summary statistics. Reference the descriptives
table generated by script.>

\input{../output/tables/tab_descriptives.tex}

\section{Empirical Strategy}
\label{sec:methods}

<Identification strategy. Estimation model. Discussion of assumptions.>

\section{Results}
\label{sec:results}

<Main findings. Reference the main results table generated by script.>

\input{../output/tables/tab_main_results.tex}

\begin{figure}[H]
\centering
\includegraphics[width=0.8\textwidth]{../output/figures/fig_main.pdf}
\caption{<Figure caption>}
\label{fig:main}
\end{figure}

\section{Robustness}
\label{sec:robustness}

<Canonical robustness checks appropriate to the design. Invoke the
\texttt{robustness-checks} skill to structure this section.>

\input{../output/tables/tab_robustness.tex}

\section{Discussion}
\label{sec:discussion}

<Main finding restated in substantive terms. Dialogue with the cited
literature: which prior debate does this resolve, refine, or complicate?
Evidence on mechanisms. External validity. Limitations.>

\section{Conclusion}
\label{sec:conclusion}

<Summary of findings. Limitations. Policy or theoretical implications.
Directions for future research.>

% --- Bibliography ---
\bibliographystyle{apalike}
\bibliography{references}
% \printbibliography  % for biblatex

% --- Appendix ---
\appendix
\section{Additional Results}
\label{app:additional}

<Supplementary tables and figures that support the main text.>

\end{document}
